% !TeX program = xelatex
\documentclass[11pt,a4paper]{article}
\usepackage[margin=25mm]{geometry}
\usepackage{amsmath,amssymb,amsthm,mathtools}
\usepackage{xcolor}
\usepackage{enumitem,microtype}
\usepackage[colorlinks=true,linkcolor=blue!50!black,citecolor=blue!50!black,urlcolor=blue!50!black]{hyperref}
\hypersetup{pdftitle={Recent Progress},pdfauthor={Santanu and Ding Haozheng}}
\setlength{\emergencystretch}{2em}
\newcommand{\C}{\mathbb C}
\newcommand{\Z}{\mathbb Z}
\newcommand{\g}{\mathfrak g}
\newcommand{\slc}{\mathfrak{sl}_2}
\newcommand{\ad}{\operatorname{ad}}
\newcommand{\im}{\operatorname{Im}}
\newcommand{\Span}{\operatorname{span}_{\C}}
\newcommand{\soc}{\operatorname{soc}}
\newcommand{\Ann}{\operatorname{Ann}}
\newcommand{\Hom}{\operatorname{Hom}}
\newcommand{\T}{\mathcal T}
\newcommand{\D}{\mathcal D}
\newcommand{\X}{\mathcal X}
\newtheorem{theorem}{Theorem}[section]
\newtheorem{proposition}[theorem]{Proposition}
\newtheorem{lemma}[theorem]{Lemma}
\newtheorem{corollary}[theorem]{Corollary}
\newtheorem{remark}[theorem]{Remark}
\title{Recent Progress\\[3pt]
  \large Boundary and Nonboundary Localization Quotients of Smooth Affine
  \texorpdfstring{\(\slc\)}{sl2}-Modules}
\author{Santanu \quad Ding Haozheng}
\date{September 26, 2026}
\begin{document}
\maketitle

\begin{abstract}
We report recent progress on localization quotients of smooth modules over
the affine Lie algebra $\g=(\slc\otimes\C[t,t^{-1}])\oplus\C K$.
The established results are as follows. (1) The boundary isomorphism
$\T_{f_b}M_{a,b}(\chi)\cong M_{a+1,b-1}(\chi^+)$ under $\lambda_L\ne0$,
together with its iteration and its compatibility with the degree
derivation. (2) For $c<b$, the quotient $Q_c=\T_{f_c}M_{a,b}(\chi)$ has a
PBW basis in the negative powers of $f_c$; $f_c$ acts surjectively and
locally nilpotently, $f_j$ acts injectively for $j\le b$, $j\ne c$, and
locally nilpotently for $j>b$; the elements $h_i$ with $1\le i\le b-c$
have no nonzero locally finite vectors, and $h_i-\lambda_i$ is locally
nilpotent for $i>b-c$; the annihilator is preserved, and quotient
localizations in distinct $f_j$ commute. (3) Under $\lambda_L\ne0$,
$\kappa\ne-2$ and $c\le-a-1$, the quotient $Q_c$ is simple if and only if
$\lambda_0-c\kappa\notin\Z$; for $\lambda_0-c\kappa\in\Z$ the subspace
$G_c=\{v:E^nv=0\text{ for some }n\}$ is the unique maximal submodule and
$Q_c/G_c$ is simple. (4) Algebraic central reduction, including the
Sugawara operators at the critical level, commutes with the quotient
localization. For the subspace $G_c$ we record the decomposition into
irreducible modules of the subalgebra generated by $E=e_{-c}$, $F=f_c$,
$H=h_0-cK$, the corrected commutation formula for the action of $E$, the
triangular system determining the vectors annihilated by $E$, and a
projection lowering the highest weight by $2$; the classification of the
submodules of $G_c$, and the simplicity in the case $-a\le c<b$, remain
open.
\end{abstract}

\section{Setup and notation}
\label{sec:setup}

Let $\g=(\slc\otimes\C[t,t^{-1}])\oplus\C K$ and $U=U(\g)$; write
$x_i=x\otimes t^i$. For $i,j\in\Z$,
\begin{equation}\label{eq:brackets}
\begin{aligned}
 [e_i,f_j]&=h_{i+j}+i\delta_{i+j,0}K,&
 [h_i,h_j]&=2i\delta_{i+j,0}K,\\
 [h_i,e_j]&=2e_{i+j},&
 [h_i,f_j]&=-2f_{i+j},
\end{aligned}
\end{equation}
$K$ is central and $[e_i,e_j]=[f_i,f_j]=0$. Fix $a,b\in\Z$ with
$a+b\ge0$, put $L=a+b+1\ge1$, and let
\[
 S_{a,b}=\Span\{K,\ h_i\ (i\ge0),\ e_i\ (i>a),\ f_j\ (j>b)\}.
\]
A character $\chi:S_{a,b}\to\C$ is given by $\kappa=\chi(K)$ and
$\lambda_i=\chi(h_i)$, $0\le i\le L$, with $\lambda_i:=0$ for $i>L$ and
$\chi=0$ on all root vectors of $S_{a,b}$. Set
\[
 M_{a,b}(\chi)=U\otimes_{U(S_{a,b})}\C_\chi,\qquad u=1\otimes1_\chi.
\]
We use the criterion of Futorny--Guo--Xue--Zhao \cite[Theorem 1.1]{FGXZ}:
\begin{equation}\label{eq:FGXZ}
 M_{a,b}(\chi)\ \text{is simple}\quad\Longleftrightarrow\quad
 \lambda_L\ne0\ \text{and}\ \kappa\ne-2.
\end{equation}

For an element $F=f_c$ we write $U_F=U[F^{-1}]$ (an Ore localization, since
$\ad F$ is locally nilpotent on $U$), $D_FV=U_F\otimes_UV$, and
\begin{equation}\label{eq:TF}
 \T_FV=\operatorname{coker}(V\to D_FV).
\end{equation}
The Poincar\'e--Birkhoff--Witt basis with $F$ placed first makes
$M_{a,b}(\chi)$ free over $\C[F]$, so $F$ acts injectively and
$\T_FM=D_FM/M$. For $c<b$ we abbreviate
\[
 Q_c=\T_{f_c}M_{a,b}(\chi),\qquad w_m=f_c^{-m}u+M\quad(m\ge1).
\]
The two cases $c\le-a-1$ (equivalently $b-c\ge L$) and $-a\le c<b$ will
be treated separately in Sections~\ref{sec:case-D} and \ref{sec:open}.

\section{The boundary isomorphism}
\label{sec:boundary}

\subsection{The basic case \texorpdfstring{$\T_{f_1}M_{-1,1}(\varphi)\cong M_{0,0}(\psi)$}{basic}}

Fix $M=M_{-1,1}(\varphi)$, $u=u_\varphi$, $F=f_1$, $T=\T_FM$,
$w_m=F^{-m}u+M$; write $\kappa=\varphi(K)$, $\lambda_i=\varphi(h_i)$, so
\begin{equation}\label{eq:cyclic}
 e_iu=0\ (i\ge0),\quad f_ju=0\ (j\ge2),\quad h_iu=\lambda_iu\ (i\ge0),
 \quad Ku=\kappa u,
\end{equation}
with $\lambda_i=0$ for $i\ge2$. Define $\psi:S_{0,0}\to\C$ by
$\psi(K)=\kappa$, $\psi(h_0)=\lambda_0+2$, $\psi(h_1)=\lambda_1$,
$\psi(h_i)=0$ ($i\ge2$), $\psi(e_i)=\psi(f_i)=0$ ($i\ge1$), and put
$A=M_{0,0}(\psi)$, $v=1\otimes1_\psi$.

For a finite nondecreasing integer sequence $\alpha=(i_1,\ldots,i_k)$ with
$i_\nu\le r$ (including the empty sequence) write
$\ell(\alpha)=k$, $x(\alpha)=x_{i_1}\cdots x_{i_k}$,
$x(\varnothing)=1$. Let
\[
 \X=\{X=e(\alpha)h(\beta)f(\gamma):\alpha,\beta\in\Gamma_{-1},\,
 \gamma\in\Gamma_0\},\qquad d(X)=2\ell(\alpha)+\ell(\beta),
\]
and set $B_{m,X}=F^{-m}Xu+M$ and
$T_{\le q}=\Span\{B_{m,X}:m\ge1,\ d(X)\le q\}$ ($T_{\le q}=0$ for
$q<0$). Then $\{F^rXu:r\ge0,\ X\in\X\}$ is a basis of $M$,
$\{F^rXu:r\in\Z,\ X\in\X\}$ is a basis of $D_FM$, and
$\{B_{m,X}:m\ge1,\ X\in\X\}$ is a basis of $T$; the source $A$ has
basis $\{Xe_0^pv:p\ge0,\ X\in\X\}$ with $e_0$ placed after all factors
of $X$ and before $S_{0,0}$. In particular $w_1\ne0$.

Put $D=\ad F$. From \eqref{eq:brackets},
\begin{equation}\label{eq:ad}
\begin{aligned}
 D(f_n)&=0,&D(h_n)&=2f_{n+1},\\
 D(e_n)&=-(h_{n+1}+n\delta_{n+1,0}K),&
 D^2(e_n)&=-2f_{n+2},\qquad D^3(e_n)=0.
\end{aligned}
\end{equation}
\begin{lemma}\label{lem:commute}
For $x\in U$ and $m\ge1$, the finite sum
\begin{equation}\label{eq:commute}
 xF^{-m}=\sum_{j\ge0}\binom{m+j-1}{j}F^{-m-j}D^j(x)
\end{equation}
holds in $U_F$.
\end{lemma}
\begin{proof}
For $m=1$, multiply the right side by $F$ on the right and use
$D^j(x)F=FD^j(x)-D^{j+1}(x)$; the sum telescopes to $x$. Suppose
\eqref{eq:commute} holds for $m$. Using
$yF^{-1}=\sum_{k\ge0}F^{-1-k}D^k(y)$ for $y=D^j(x)$ gives
\[
 xF^{-(m+1)}=\sum_{j,k\ge0}\binom{m+j-1}{j}F^{-m-j-1-k}D^{j+k}(x),
\]
whose coefficient of $F^{-(m+1)-r}D^r(x)$ is
$\sum_{j=0}^{r}\binom{m+j-1}{j}=\binom{m+r}{r}$, as required.
\end{proof}

Combining \eqref{eq:commute} with \eqref{eq:ad} gives the commutation
formulas
\begin{equation}\label{eq:actions}
\begin{aligned}
 f_nF^{-m}&=F^{-m}f_n,\\
 h_nF^{-m}&=F^{-m}h_n+2mF^{-m-1}f_{n+1},\\
 e_nF^{-m}&=F^{-m}e_n-mF^{-m-1}(h_{n+1}+n\delta_{n+1,0}K)
 -m(m+1)F^{-m-2}f_{n+2}.
\end{aligned}
\end{equation}
Applying them to $u$ with \eqref{eq:cyclic} and $m=1$ gives
\begin{equation}\label{eq:w1}
\begin{gathered}
 e_nw_1=f_nw_1=0\ (n\ge1),\qquad Kw_1=\kappa w_1,\\
 h_0w_1=(\lambda_0+2)w_1,\qquad h_1w_1=\lambda_1w_1,\qquad
 h_nw_1=0\ (n\ge2),
\end{gathered}
\end{equation}
because $f_1w_1=u+M=0$; the corrections for $e_n$ vanish by
$h_{n+1}u=f_{n+2}u=0$ ($n\ge1$), those for $h_n$ by $f_{n+1}u=0$
($n\ge1$), and the correction for $h_0$ is $2F^{-2}f_1u+M=2w_1$. Thus
$sw_1=\psi(s)w_1$ for all $s\in S_{0,0}$; since the unital algebra
character $\psi_U:U(S_{0,0})\to\C$ extends $\psi$, we get
$aw_1=\psi_U(a)w_1$ for $a\in U(S_{0,0})$. Hence the map
\[
 \Phi:A\longrightarrow T,\qquad x\otimes z1_\psi\longmapsto z\,xw_1,
\]
kills every balancing relation
$x a\otimes z1_\psi-x\otimes a(z1_\psi)$, and descends to a $U$-linear
$\Phi$. From \eqref{eq:actions} with $n=0$,
$e_0w_m=-m\lambda_1w_{m+1}$, hence
\begin{equation}\label{eq:tower}
 e_0^pw_1=(-1)^pp!\lambda_1^pw_{p+1}\quad(p\ge0).
\end{equation}
\begin{lemma}\label{lem:lower}
For $m\ge1$ and $X\in\X$,
\[
 Xw_m=B_{m,X}+R_{m,X},\qquad R_{m,X}\in T_{\le d(X)-1}.
\]
\end{lemma}
\begin{proof}
Give the generators weights $\mathrm{wt}(e_i)=2$, $\mathrm{wt}(h_i)=1$,
$\mathrm{wt}(f_i)=\mathrm{wt}(K)=0$, and filter $U$ by total word weight,
$U_{\le q}=0$ for $q<0$. By \eqref{eq:brackets} the weight of a bracket is
at most the sum of the input weights: for $[h,e]$, $[e,f]$, $[h,f]$, $[h,h]$
the output weights are at most $2,1,0,0$ against input sums $3,2,1,2$. So
reordering in the PBW basis never increases the weight, and neither does
evaluation of the generators of $S_{0,0}$ on $u$. By \eqref{eq:ad} and the
Leibniz rule, $D(U_{\le q})\subseteq U_{\le q-1}$, hence
$D^j(X)\in U_{\le d(X)-j}$ and $D^j(X)=0$ for $j>d(X)$. Expanding in the
PBW basis of $M$,
\[
 D^j(X)u=\sum_{r\ge0,\,X'\in\X}c_{j,r,X'}F^rX'u,\qquad
 c_{j,r,X'}\ne0\Rightarrow d(X')\le d(X)-j.
\]
Substitute into \eqref{eq:commute}. The $j=0$ term is exactly $B_{m,X}$.
For $j\ge1$, a term with $r\ge m+j$ lies in $M$ and vanishes in $T$; every
surviving term is $B_{m+j-r,X'}$ with $d(X')\le d(X)-j<d(X)$.
\end{proof}

\begin{theorem}\label{thm:basic}
If $\lambda_1\ne0$, then $\Phi:A\xrightarrow{\ \sim\ }T$, and $\Phi$ is
the unique $\g$-module isomorphism sending $v$ to $w_1$. This holds for
every $\kappa\in\C$.
\end{theorem}
\begin{proof}
\emph{Surjectivity.} Prove by induction on $d\ge0$ the statement
$P(d)$: $B_{m,X}\in\im\Phi$ for every $m\ge1$ and every $X\in\X$ with
$d(X)\le d$ (all negative powers simultaneously). For $d=0$ we have
$X=f(\gamma)$, which commutes with $F$, so $B_{m,X}=Xw_m=\Phi(Xq_m)$ with
$q_m=\frac{(-1)^{m-1}}{(m-1)!\lambda_1^{m-1}}e_0^{m-1}v$, using
\eqref{eq:tower}. If $d\ge1$ and $P(d-1)$ holds, Lemma~\ref{lem:lower}
gives $\Phi(Xq_m)=B_{m,X}+\sum_\nu c_\nu B_{m_\nu,X_\nu}$ with
$d(X_\nu)\le d-1$; by induction choose $a_\nu$ with
$\Phi(a_\nu)=B_{m_\nu,X_\nu}$, and
$B_{m,X}=\Phi(Xq_m-\sum_\nu c_\nu a_\nu)$. The recursion strictly
decreases the nonnegative integer $d$ and terminates. For example
$d(h_{-1})=1$ and $h_{-1}w_m=B_{m,h_{-1}}+2mB_{m+1,f_0}$, giving the
preimage $B_{m,h_{-1}}=\Phi(h_{-1}q_m-2mf_0q_{m+1})$.

\emph{Injectivity.} By \eqref{eq:tower} and Lemma~\ref{lem:lower},
\begin{equation}\label{eq:leading}
 \Phi(Xe_0^pv)\equiv(-1)^pp!\lambda_1^pB_{p+1,X}
 \pmod{T_{\le d(X)-1}}.
\end{equation}
Let $0\ne a=\sum_{X,p}a_{X,p}Xe_0^pv\in A$ and
$d=\max\{d(X):a_{X,p}\ne0\}$. In $T_{\le d}/T_{\le d-1}$,
\[
 \Phi(a)+T_{\le d-1}
 =\sum_{\substack{X,p\\d(X)=d}}a_{X,p}(-1)^pp!\lambda_1^p
 (B_{p+1,X}+T_{\le d-1}),
\]
and these cosets are linearly independent; since $(-1)^pp!\lambda_1^p\ne0$,
$\Phi(a)\ne0$.
\end{proof}

\subsection{The general boundary and iteration}
\begin{theorem}\label{thm:N}
Let $N\ge1$, $M=M_{-1,N}(\varphi)$, $F=f_N$, $w_m=F^{-m}u+M$, and assume
$\lambda_N\ne0$. Define $\psi:S_{0,N-1}\to\C$ by $\psi(K)=\kappa$,
$\psi(h_i)=\lambda_i+2\delta_{i,0}$ ($i\ge0$), $\psi(e_i)=0$ ($i\ge1$),
$\psi(f_j)=0$ ($j\ge N$). Then
\[
 \Phi_N:M_{0,N-1}(\psi)\xrightarrow{\ \sim\ }\T_{f_N}M,
 \qquad v\longmapsto f_N^{-1}u+M.
\]
\end{theorem}
\begin{proof}
In \eqref{eq:ad}--\eqref{eq:actions} replace $n+1,n+2$ by $n+N,n+2N$.
Since $e_iu=0$ ($i\ge0$), $f_ju=0$ ($j>N$) and $\lambda_i=0$ ($i>N$), we get
$e_iw_1=0$ ($i\ge1$), $f_jw_1=0$ ($j\ge N$) (with $f_Nw_1=u+M=0$) and
$h_iw_1=(\lambda_i+2\delta_{i,0})w_1$, so the balancing argument gives a
well-defined $\Phi_N$. As $e_0u=f_{2N}u=0$ and $h_Nu=\lambda_Nu$,
\[
 e_0w_m=-m\lambda_Nw_{m+1},\qquad
 w_m=\Phi_N\Bigl(\frac{(-1)^{m-1}e_0^{m-1}v}{(m-1)!\lambda_N^{m-1}}\Bigr).
\]
Replace $\X$ by
$\X_N=\{e(\alpha)h(\beta)f(\gamma):\alpha,\beta\in\Gamma_{-1},\,
\gamma\in\Gamma_{N-1}\}$; the source and target bases are
$\{Xe_0^pv\}$ and $\{F^{-m}Xu+M\}$, and the same degree
$d(X)=2\ell(\alpha)+\ell(\beta)$ is strictly lowered by $\ad f_N$ while
reordering and the relations in the inducing subalgebra do not increase it.
The same strong induction proves surjectivity, and the coefficients
$(-1)^pp!\lambda_N^p\ne0$ prove injectivity.
\end{proof}

\begin{corollary}\label{cor:general}
Let $a,b\in\Z$, $a+b\ge0$, $L=a+b+1$, $M=M_{a,b}(\chi)$. If
$\lambda_L\ne0$, then
\[
 M_{a+1,b-1}(\chi^+)\xrightarrow{\ \sim\ }\T_{f_b}M,
 \qquad v\longmapsto f_b^{-1}u+M,
\]
where $\chi^+$ vanishes on all root vectors of $S_{a+1,b-1}$ and
$\chi^+(K)=\kappa$, $\chi^+(h_i)=\lambda_i+2\delta_{i,0}$.
\end{corollary}
\begin{proof}
Set $F=f_b$, $Y=e_{a+1}$, $w_m=F^{-m}u+M$. For $i\ge a+2$, commuting
$e_i$ past $F^{-1}$ produces the Cartan index $i+b\ge L+1$ and the central
$\delta$-term, which vanish on $u$, and the remaining correction vanishes
because $i+2b>b$, so $f_{i+2b}u=0$. The other inducing relations are
$f_jw_1=0$ ($j\ge b$), $h_iw_1=(\lambda_i+2\delta_{i,0})w_1$ and
$Kw_1=\kappa w_1$. Since $Yu=0$, $L>0$ and $f_{a+1+2b}u=0$,
\[
 Yw_m=-m\lambda_Lw_{m+1}.
\]
Use $\X_{a,b}=\{e(\alpha)h(\beta)f(\gamma):\alpha\in\Gamma_a,\,
\beta\in\Gamma_{-1},\ \gamma\in\Gamma_{b-1}\}$ and the source basis
$\{XY^pv\}$. The derivation $\ad f_b$ strictly lowers
$d(X)=2\ell(\alpha)+\ell(\beta)$, so the same induction and highest-degree
comparison apply with coefficients $(-1)^pp!\lambda_L^p\ne0$.
\end{proof}

Iterating Corollary~\ref{cor:general} (applied at successive boundaries,
each step requiring the corresponding $\lambda$ to be nonzero) yields
$\T_{f_c}\cdots\T_{f_b}M_{a,b}\cong M_{a+b-c+1,c-1}(\chi^{(b-c+1)})$
with $\chi^{(k)}(h_0)=\lambda_0+2k$ and all other parameters fixed.

\subsection{The derivation}
Let $\widetilde{\g}=\g\rtimes\C d$ with $[d,x_i]=ix_i$, and
$\mathcal I(V)=U(\widetilde{\g})\otimes_UV$. Then
$\mathcal I(V)\cong\C[d]\otimes V$ with
\begin{equation}\label{eq:derivation-action}
 x_i(P(d)\otimes v)=P(d-i)\otimes x_iv.
\end{equation}
\begin{proposition}\label{prop:derivation}
For $F=f_c$ acting injectively on $V$,
\[
 \T_F\mathcal I(V)\cong\mathcal I(\T_FV),
\qquad
 F^{-m}(P(d)\otimes v)\longmapsto P(d+mc)\otimes F^{-m}v.
\]
\end{proposition}
\begin{proof}
$\mathcal I$ is exact and faithful. It remains to check that the displayed
map intertwines the action. For $x=h_r$:
\[
 h_rF^{-m}P(d)\otimes v
 =F^{-m}h_rP(d)\otimes v+2mF^{-m-1}f_{r+c}P(d)\otimes v
 =F^{-m}P(d-r)\otimes h_rv+2mF^{-m-1}P(d-(r+c))\otimes f_{r+c}v,
\]
which is sent to
$P(d-r+mc)\otimes F^{-m}h_rv+2mP(d-(r+c)+c(m+1))\otimes F^{-m-1}f_{r+c}v
=P(d+mc-r)\otimes h_rF^{-m}v$, the image of the right side. For $x=e_r$
the same computation uses
\[
 e_rF^{-m}=F^{-m}e_r-mF^{-m-1}(h_{r+c}+r\delta_{r+c,0}K)
 -m(m+1)F^{-m-2}f_{r+2c},
\]
and $f_r$ commutes with $F$.
\end{proof}

\section{Structure of the quotients \texorpdfstring{$Q_c$}{Qc}, \texorpdfstring{$c<b$}{c<b}}
\label{sec:nonboundary}

Throughout this section $c<b$. Let $\X$ be the set of ordered monomials
$X=f_{j_1}\cdots f_{j_p}e_{i_1}\cdots e_{i_q}h_{r_1}\cdots h_{r_s}$ with
$j_t\le b$, $j_t\ne c$, $i_t\le a$, $r_u<0$ (including the empty word), and
put $\deg X=2q-2p$. For a nonzero $v\in Q_c$ we denote by $n_F(v)$ the
smallest integer $m\ge1$ such that $F^mv=0$; in terms of the expansion
$v=\sum_{m\ge1,X}c_{m,X}F^{-m}Xu+M$, $n_F(v)$ is the largest $m$ with
$c_{m,X}\ne0$ for some $X$.

\begin{proposition}\label{prop:pbw}
The set $\{F^{-m}Xu+M:m\ge1,\ X\in\X\}$ is a basis of $Q_c$, and
\[
 \ker(F^r|_{Q_c})=\Span\{F^{-m}Xu+M:1\le m\le r,\ X\in\X\}.
\]
The elements $f_j$ act as follows: $f_c$ surjectively and locally
nilpotently; $f_j$ injectively but not surjectively for $j\le b$, $j\ne c$;
$f_j$ locally nilpotently for $j>b$. Moreover $Q_c$ is a smooth weight
module with $K=\kappa$ and $h_0$ semisimple,
$\operatorname{Supp}_{h_0}Q_c=\lambda_0+2\Z$, each nonzero weight space
has countable dimension, and $Q_c$ has no nonzero finite-dimensional
submodules or quotients.
\end{proposition}
\begin{proof}
The basis and the kernel formula follow from the PBW basis with $F$ placed
first and from $M$ being free over $\C[F]$. For the elements $f_j$:
$F^r(F^{-m}Xu)=F^{-m+r}Xu$ vanishes for $r\ge m$ and
$F^{-m}Xu=F(F^{-m-1}Xu)$, giving surjectivity and local nilpotence of
$f_c$. For $c<j\le b$, the element $f_j$ is a PBW basis element of the
complement of $S_{a,b}$ commuting with $F$, so $f_j(F^{-m}Xu)=F^{-m}f_jXu$
is multiplication by an independent variable; it is injective, and its
image consists of vectors whose expansion is divisible by $f_j$, so $w_1$
has no preimage and $f_j$ is not surjective. For $j>b$: $f_ju=0$, $f_j$
commutes with $F$ and with all $f_k$, and the adjoint action of $f_j$
raises indices, since $[f_j,h_r]=2f_{j+r}$ with $j+r>b$ and
$[f_j,e_i]=-h_{j+i}-j\delta_{j+i,0}K$; hence iterating $f_j$ on a PBW
basis vector eventually produces only elements that annihilate $u$, and
$f_j$ is locally nilpotent. The weight formula is
$h_0(F^{-m}Xu)=(\lambda_0+2m+\deg X)F^{-m}Xu$, and powers of $h_{-1}$
produce infinitely many independent vectors in each weight space.
Smoothness follows from \eqref{eq:comm-c} below: for a fixed PBW basis
vector, $r,r+c,r+2c$ large enough annihilate $u$ and the central
$\delta$-term vanishes. A finite-dimensional quotient cannot carry $f_c$
both surjective and locally nilpotent, and a finite-dimensional smooth
module is trivial, which is excluded because $f_{c+1}$ is injective; this
proves the last assertion.
\end{proof}

In the full localization, for all $i\in\Z$ and $m\ge1$,
\begin{equation}\label{eq:comm-c}
\begin{aligned}
 h_iF^{-m}&=F^{-m}h_i+2mF^{-m-1}f_{i+c},\\
 e_iF^{-m}&=F^{-m}e_i-mF^{-m-1}(h_{i+c}+i\delta_{i+c,0}K)
 -m(m+1)F^{-m-2}f_{i+2c},
\end{aligned}
\end{equation}
obtained from \eqref{eq:commute} with $D=\ad f_c$:
$D(f_n)=0$, $D(h_n)=2f_{n+c}$, $D(e_n)=-(h_{n+c}+n\delta_{n+c,0}K)$,
$D^2(e_n)=-2f_{n+2c}$, $D^3(e_n)=0$.

\begin{proposition}\label{prop:cartan}
If $1\le i\le b-c$, then for $0\ne v\in Q_c$, $z\in\C$ and
$0\ne p\in\C[t]$,
\[
 n_F((h_i-z)v)=n_F(v)+1,\qquad p(h_i)v\ne0.
\]
If $i>b-c$, then $h_i-\lambda_i$ is locally nilpotent on $Q_c$. Hence
$Q_c$ is not isomorphic to any standard module $M_{a',b'}(\psi)$, and
\[
 \Hom_\g(M_{a',b'}(\psi),Q_c)=0,\qquad
 \Hom_\g(\T_{f_c}M,\T_{f_{c'}}M)=0\quad(c\ne c',\ c,c'\le b).
\]
\end{proposition}
\begin{proof}
Write $v=\sum_{m,X}c_{m,X}F^{-m}Xu$ and let $m_0=n_F(v)$. By
\eqref{eq:comm-c},
\[
 h_iv=\sum_{m,X}c_{m,X}\bigl(F^{-m}h_iXu+2mF^{-m-1}f_{i+c}Xu\bigr).
\]
Applying $F^{m_0+1}$ kills every term, so $F^{m_0+1}(h_i-z)v=0$. On the
other hand, applying $F^{m_0}$ kills all terms except the terms with
$m=m_0$ of the second sum, which contribute
$2m_0F^{-1}f_{i+c}X_{m_0}u$; here $c+i\le b$ and $c+i\ne c$, so $f_{i+c}$
is a PBW basis element of the complement of $S_{a,b}$, and this vector is
nonzero in $Q_c$. Thus $n_F((h_i-z)v)=m_0+1$. Iterating,
$n_F(p(h_i)v)=m_0+\deg p$, so $p(h_i)v\ne0$.

If $i>b-c$, then $i+c>b$, so $h_iw_m=\lambda_iw_m$, and for a PBW basis
vector, $(h_i-\lambda_i)^NXw_m=(\ad h_i)^N(X)w_m$. By the Leibniz rule,
each term of $(\ad h_i)^N(X)$ is a product in which each factor is either
unchanged, or has received $k\ge1$ adjoint actions of $h_i$ and has index
increased by $ki$, or is a scalar multiple of $K$ arising from
$[h_i,h_{-i}]$. An $f$-factor can receive at most $\lfloor(b-j)/i\rfloor$
adjoint actions before its index exceeds $b$ and it annihilates $u$; an
$e$-factor at most $\lfloor(a-j)/i\rfloor$; an $h$-factor at most one
(because $[h_i,[h_i,h_{-i}]]=0$ and $[h_i,K]=0$). Hence for
$N>\sum_{t}\lfloor(b-j_t)/i\rfloor+\sum_{t}\lfloor(a-j_t)/i\rfloor+s$
every term of $(\ad h_i)^N(X)$ either contains a factor annihilating $u$
or vanishes, so $(h_i-\lambda_i)^NXw_m=0$.

A standard module has a cyclic vector annihilated by the root vectors of
its inducing subalgebra and which is an eigenvector of $h_1$; the image of
such a vector under a homomorphism into $Q_c$ would be a nonzero vector on
which $h_1$ acts by a scalar, contradicting $n_F((h_1-\lambda)v)=n_F(v)+1$
for $\lambda\in\C$. This gives the first Hom-vanishing. The second
follows since $f_c$ is locally nilpotent on $\T_{f_c}M$ but injective on
$\T_{f_{c'}}M$ for $c'\ne c$ (Proposition~\ref{prop:pbw}).
\end{proof}

\begin{proposition}[Commuting quotient localizations]
\label{prop:commuting}
For $j\le b$, $j\ne c$,
\[
 \T_{f_j}\T_{f_c}M\cong\frac{D_{f_c,f_j}M}{D_{f_c}M+D_{f_j}M}
 \cong\T_{f_c}\T_{f_j}M.
\]
\end{proposition}
\begin{proof}
From $0\to M\to D_{f_c}M\to\T_{f_c}M\to0$ and the flatness of the Ore
localization at $f_j$,
$D_{f_j}\T_{f_c}M\cong D_{f_c,f_j}M/D_{f_j}M$. Quotienting by the image of
$\T_{f_c}M$ corresponds to dividing by $(D_{f_c}M+D_{f_j}M)/D_{f_j}M$;
the displayed quotient is symmetric in $f_c,f_j$.
\end{proof}

\begin{proposition}[Annihilators, nonsplitting, derived functors]
\label{prop:ann}
In the ordinary enveloping algebra,
\[
 \Ann_UM=\Ann_UD_FM=\Ann_UQ_c,
\]
the sequence $0\to M\to D_FM\to Q_c\to0$ does not split,
$\Hom_\g(Q_c,M)=\Hom_\g(Q_c,D_FM)=0$, and for
$\T_FV=(U_F/U)\otimes_UV$ one has
\[
 L_1\T_F(V)\cong\{v\in V:F^nv=0\text{ for some }n\ge1\},\qquad
 L_j\T_F(V)=0\ (j\ge2).
\]
In particular $D_FQ_c=\T_FQ_c=0$ and $L_1\T_F(Q_c)\cong Q_c$.
\end{proposition}
\begin{proof}
$\Ann_UM$ is a two-sided ideal, hence stable under $\Delta=\ad F$; the
commutation formula \eqref{eq:commute} then shows that it annihilates
$D_FM$, and $M\hookrightarrow D_FM$ gives
$\Ann_UD_FM\subseteq\Ann_UM\subseteq\Ann_UQ_c$. Conversely, if
$sQ_c=0$ and $\Delta^{t+1}(s)=0$, then for $v\in M$ and $m\ge1$ we have
$sF^{-m}v\in M$, and multiplying \eqref{eq:commute} by $F^{m+t}$ gives
\[
 P_v(m):=\sum_{j=0}^{t}\binom{m+j-1}{j}F^{t-j}\Delta^j(s)v\in F^{m+t}M.
\]
This is a vector-valued polynomial in $m$ whose coefficients involve only
finitely many PBW basis vectors, so their nonnegative $F$-degrees have a
common upper bound; for $m$ large the span meets $F^{m+t}M$ trivially.
Hence $P_v(m)=0$ for all large $m$, so the polynomial is identically zero;
at $m=0$ only $j=0$ survives ($\binom{j-1}{j}=0$), giving $F^tsv=0$, and $F$
is injective on $M$. The splitting and Hom-vanishing follow from the
$F$-torsion of $Q_c$ versus the injectivity of $F$ on $M$ and $D_FM$. For
the derived functors use the flat resolution
$0\to U\to U_F\to U_F/U\to0$ of the right $U$-module $U_F/U$:
$\T_FV=(U_F/U)\otimes_UV$ and $L_1\T_FV=\ker(V\to D_FV)$ is the
$F$-torsion, while $L_j=0$ for $j\ge2$ by length one.
\end{proof}

\section{The case \texorpdfstring{$c\le-a-1$}{c<=-a-1}: simplicity and the unique maximal submodule}
\label{sec:case-D}

Assume $c\le-a-1$ and set
\begin{equation}\label{eq:EHF}
 E=e_{-c},\qquad F=f_c,\qquad H=h_0-cK,\qquad \mu=\lambda_0-c\kappa.
\end{equation}
Then $[H,E]=2E$, $[H,F]=-2F$, $[E,F]=H$, so the subalgebra generated by
$E,H,F$ is isomorphic to $\slc$; the assumption $c\le-a-1$ means $-c>a$,
so $Eu=0$ and $Hu=\mu u$, and $E$ is locally nilpotent on $M$. The identity
\begin{equation}\label{eq:triple}
 EF^{-m}=F^{-m}E-mF^{-m-1}H-m(m+1)F^{-m-2}F
\end{equation}
gives, on $w_m=F^{-m}u+M$,
\begin{equation}\label{eq:pure}
 Ew_m=-m(\mu+m+1)w_{m+1},\qquad Hw_m=(\mu+2m)w_m.
\end{equation}

Define
\[
 G_c=\{v\in Q_c:E^nv=0\text{ for some }n\ge1\}.
\]
Since $\ad E$ is locally nilpotent, $G_c$ is a $\g$-submodule: for
$v\in G_c$ and $s\in U$,
$E^r(sv)=\sum_{j=0}^{r}\binom rj(\ad E)^j(s)E^{r-j}v$, which vanishes for
$r$ large. Every nonzero submodule of $G_c$ contains a nonzero vector of
$P_\nu:=\ker E\cap G_c\cap(Q_c)_\nu$ for some $\nu$: take the last
nonzero $E$-image.

\begin{proposition}\label{prop:G}
Without assuming \eqref{eq:FGXZ}, $G_c$ is a proper submodule of $Q_c$ and
\[
 G_c=0\ \Longleftrightarrow\ \mu\notin\Z.
\]
If $\mu\in\Z$ and $k\ge\max(1,\mu+2)$, then
$0\ne f_b^kw_1\in G_c$ and $E^{2k-\mu-1}(f_b^kw_1)=0$. Moreover $Q_c$ is
cyclic, $Q_c=Uw_1$ if and only if $\mu\notin\{-2,-3,\ldots\}$, and for
$\mu=-s$, $s\ge2$, one has $Q_c=Uw_s\ne Uw_1$. The operator $E$ is
injective on $Q_c/G_c$.
\end{proposition}
\begin{proof}
By \eqref{eq:pure}, $w_m\notin G_c$ for $m$ large (all coefficients
$-m(\mu+m+1)$ are nonzero from some point on), so $G_c$ is proper. If
$0\ne v\in\ker E$ is an $H$-weight vector of weight $\nu$ and $n$ is
minimal with $F^nv=0$, then
\[
 0=EF^nv=\sum_{i=0}^{n-1}F^i[E,F]F^{n-1-i}v
 =\sum_{i=0}^{n-1}(\nu-2(n-1-i))F^{n-1}v
 =n(\nu-n+1)F^{n-1}v,
\]
so $\nu=n-1\in\Z$; but $\nu\in\mu+2\Z$, whence $G_c=0$ for
$\mu\notin\Z$ (every nonzero vector of $G_c$ has a nonzero weight
component, and its last nonzero $E$-image lies in $\ker E$).

For $\mu\in\Z$, put $Y=f_b$ and $\eta=\lambda_{b-c}$. Since
$[E,Y]=h_{b-c}$ (the $\delta$-term vanishes because $b-c\ge1$) and
$[H,Y]=-2Y$, while $[h_{b-c},Y]=-2f_{b+(b-c)}$ annihilates $u$, we get the
commutation formula on $\C[F,Y]u$:
\begin{equation}\label{eq:two-var}
 E(F^jY^\ell u)=j(\mu-j+1-2\ell)F^{j-1}Y^\ell u
 +\ell\eta F^jY^{\ell-1}u.
\end{equation}
Indeed $EF^j=F^jE+jF^{j-1}(H-j+1)$ and
$[E,Y^\ell]u=\ell Y^{\ell-1}h_{b-c}u=\ell\eta Y^{\ell-1}u$. Define
$v_k=\sum_{j=0}^{k}a_jF^jY^{k-j}u$ with
\begin{equation}\label{eq:vk}
 a_0=1,\qquad a_j=-\frac{(k-j+1)\eta}{j(\mu-2k+j+1)}a_{j-1}.
\end{equation}
The denominators are nonzero because $\mu-2k+j+1\le\mu-k+1<0$. In
\eqref{eq:two-var} the coefficient of $F^{j-1}Y^{k-j}u$ in $Ev_k$ is
$j(\mu-2k+j+1)a_j+(k-j+1)\eta a_{j-1}=0$, so $Ev_k=0$ and
$Hv_k=(\mu-2k)v_k$. Set $\nu=\mu-2k\le-2$. By \eqref{eq:triple},
$E(F^{-m}v_k)=-m(\nu+m+1)F^{-m-1}v_k$ for all $m\ge1$; iterating,
\[
 E^r(F^{-1}v_k)=\Bigl(\prod_{t=1}^{r}(-t)(\nu+t+1)\Bigr)F^{-1-r}v_k,
\]
and the factor with $t=-\nu-1$ vanishes, so
$E^{-\nu-1}(F^{-1}v_k)=0$. Passing to the quotient,
$F^{-1}v_k+M=Y^kw_1\ne0$ (the terms with $j\ge1$ land in $M$), giving
$0\ne Y^kw_1\in G_c$ and $E^{2k-\mu-1}(Y^kw_1)=0$.

The module is cyclic: \eqref{eq:pure} produces $w_{m+1}$ from $w_m$ whenever
$-m(\mu+m+1)\ne0$, and $Fw_m=w_{m-1}$ ($w_0=0$) produces the lower ones;
to pass from the $w_m$ to arbitrary PBW basis vectors, grade the words by
the weight $\rho=2\#e+\#h$: commuting $X$ past $F^{-m}$ gives
$Xw_m=F^{-m}Xu+M+$(terms of strictly lower $\rho$), and induction on
$\rho$, simultaneously for all $m$, proves that the vectors $w_m$ generate
$Q_c$. For $\mu=-s$, $s\ge2$, the coefficient in \eqref{eq:pure}
vanishes at $m=s-1$, so $w_1\in G_c$ and $Uw_1\subseteq G_c\ne Q_c$,
while all coefficients from $w_s$ upward are nonzero, so $Q_c=Uw_s$.
Finally, if $Ev\in G_c$ then $v\in G_c$, so $E$ is injective on
$Q_c/G_c$.
\end{proof}

\subsection{The containment of all proper submodules}
\label{sec:abstract}

We isolate the abstract mechanism. Let $\mathfrak a$ be a complex Lie
algebra containing elements $E,H,F$ with
$[H,E]=2E$, $[H,F]=-2F$, $[E,F]=H$; write $A=U(\mathfrak a)$ and assume:
\begin{enumerate}[label=(\roman*)]
\item $\ad E$ and $\ad F$ are locally nilpotent on $A$, and the Ore
localization $A_F$ at the powers of $F$ exists;
\item $M$ is a simple $A$-module on which $H$ acts semisimply, $E$ acts
locally nilpotently, and $F$ acts injectively.
\end{enumerate}
Let $D=A_F\otimes_AM$, $Q=D/M$, and
$G=\{q\in Q:E^nq=0\text{ for some }n\ge1\}$. The $H$-action on $D$ and
$Q$ is semisimple since $[H,F^{-m}]=2mF^{-m}$, and $G$ is an $A$-submodule
by the Leibniz expansion. Define the Casimir element of the subalgebra
generated by $E,H,F$:
\begin{equation}\label{eq:C}
 C=H^2+2H+4FE.
\end{equation}
\begin{lemma}\label{lem:casimir}
$C$ commutes with $E,H,F$. For $Hz=\nu z$ and $n\ge0$,
\begin{equation}\label{eq:factorization}
 F^nE^nz=A_n(C;\nu)z,\qquad
 A_n(t;\nu)=4^{-n}\prod_{j=0}^{n-1}
 \bigl(t-(\nu+2j)(\nu+2j+2)\bigr),
\end{equation}
with $A_0=1$. Consequently $C$ acts locally finitely on $M$ and on $D$.
\end{lemma}
\begin{proof}
$[C,E]=[H^2,E]+2[H,E]+4[FE,E]=4EH+4E+4E-4HE=0$, since
$[H^2,E]=H[H,E]+[H,E]H=2HE+2EH=4EH+4E$ and $[F,E]=-H$; similarly
$[C,F]=0$. For $n=1$,
$FEz=\tfrac14(C-H^2-2H)z=\tfrac14(C-(\nu^2+2\nu))z=A_1(C;\nu)z$.
If \eqref{eq:factorization} holds for $n$, apply it to the weight vector
$Ez$ (weight $\nu+2$) and multiply by $F$ on the left:
$F^{n+1}E^{n+1}z=F\cdot A_n(C;\nu+2)Ez=A_n(C;\nu+2)FEz
=A_{n+1}(C;\nu)z$. If $z\in M$ is a weight vector with $E^nz=0$, then
$A_n(C;\nu)z=F^nE^nz=0$ is a nonzero polynomial in $C$ annihilating $z$;
each vector of $M$ has finitely many weight components, so a product of
such polynomials kills it. As $C$ commutes with $F^{-1}$, the same
polynomials work in $D$.
\end{proof}

For an $A$-submodule $P\subseteq D$ define
\begin{equation}\label{eq:W}
 W(P)=\bigcap_{n\ge0}F^nP.
\end{equation}
\begin{lemma}\label{lem:W}
$W(P)$ is an $A$-submodule of $D$, is stable under $F^{-1}$, and if
$W(P)\ne0$ then $P=D$.
\end{lemma}
\begin{proof}
For $s\in A$ choose $\ell$ with $(\ad F)^{\ell+1}(s)=0$. We use the
binomial identity
\begin{equation}\label{eq:binomial}
 sF^{n+\ell}=\sum_{j=0}^{\ell}(-1)^j\binom{n+\ell}{j}
 F^{n+\ell-j}(\ad F)^j(s),
\end{equation}
which holds in $A$: for $N=1$ it reads $sF=Fs-[F,s]$, and the induction
step is
$sF^{N+1}=\sum_j(-1)^j\binom Nj F^{N-j}(\ad F)^j(s)F
=\sum_j(-1)^j\binom NjF^{N-j+1}(\ad F)^j(s)
+\sum_{j\ge1}(-1)^j\binom N{j-1}F^{N-j+1}(\ad F)^j(s)$,
whose coefficients $\binom Nj+\binom N{j-1}=\binom{N+1}{j}$ give the
identity for $N+1$. If $w\in W(P)$, write $w=F^{n+\ell}p$ with $p\in P$;
then $sw=\sum_j(-1)^j\binom{n+\ell}{j}F^{n+\ell-j}(\ad F)^j(s)p\in F^nP$
for every $n$, so $sw\in W(P)$. Writing $w=F^{n+1}p$ gives
$F^{-1}w=F^np\in W(P)$, so $W(P)$ is $F^{-1}$-stable. If $0\ne w\in W(P)$,
write $w=F^{-m}z$ with $z\in M$; then $F^mw=z\in M\cap W(P)$ is nonzero
($F$ is invertible on $D$). Since $M$ is simple, $M\subseteq W(P)$, and
$F^{-1}$-stability gives $D=A_FM\subseteq W(P)\subseteq P$, so $P=D$.
\end{proof}

\begin{theorem}[Containment of all proper submodules]
\label{thm:containment}
Under (i)--(ii), every proper $A$-submodule $N\subsetneq Q$ is contained in
$G$. If $G\ne Q$, then $G$ is the unique maximal submodule of $Q$, $Q/G$ is
simple, $Q$ is indecomposable, and if $G\ne0$ the sequence
$0\to G\to Q\to Q/G\to0$ does not split.
\end{theorem}
\begin{proof}
Let $N\subseteq Q$ be a submodule not contained in $G$, and let
$P\subseteq D$ be its inverse image, so $M\subseteq P$ and $P/M=N$. Choose
$q\in N$ on which $E$ is not locally nilpotent; since $H$ is semisimple and
each vector has finitely many weight components, one weight component $q_0$
of $q$ has the same property, and $q_0\in N$ (it is obtained from $q$ by a
polynomial in $H$). Lift $q_0$ to a weight vector $w\in P$, $Hw=\eta w$. By
Lemma~\ref{lem:casimir} choose $0\ne p(t)$ with $p(C)w=0$. The set of
nonnegative integers $k$ with $p((\eta+2k)(\eta+2k+2))=0$ is finite, each
root of $p$ contributing at most two values of $k$; choose $r$ larger than
all of them and set $v=E^rw$, $\nu=\eta+2r$. Then $E^rq_0\ne0$, so
$v\ne0$; also $Hv=\nu v$, $p(C)v=0$ and
$p((\nu+2j)(\nu+2j+2))\ne0$ for all $j\ge0$.

For each $n\ge1$, the polynomials $p(t)$ and $A_n(t;\nu)$ are therefore
coprime, so there are $B_n,R_n\in\C[t]$ with
\[
 A_n(t;\nu)B_n(t)+p(t)R_n(t)=1.
\]
The vector $B_n(C)v$ has weight $\nu$, and
\eqref{eq:factorization} gives
\[
 F^n\bigl(E^nB_n(C)v\bigr)=A_n(C;\nu)B_n(C)v=v,
\]
so $v\in F^nP$ for every $n\ge0$: $0\ne v\in W(P)$. Lemma~\ref{lem:W}
gives $P=D$, hence $N=Q$. The remaining assertions follow: if $G\ne Q$ it
is a proper submodule containing every proper submodule; a nontrivial
direct sum decomposition of $Q$ would put both summands in $G$; a splitting
of the displayed sequence would give such a decomposition when $G\ne0$.
\end{proof}

\begin{theorem}[Simplicity criterion and the simple quotient]
\label{thm:deep}
Assume
\begin{equation}\label{eq:deep-assumptions}
 \lambda_L\ne0,\qquad \kappa\ne-2,\qquad c<b,\qquad c\le-a-1.
\end{equation}
Then
\[
 \boxed{\quad Q_c\ \text{is simple}\quad\Longleftrightarrow\quad
 \mu=\lambda_0-c\kappa\notin\Z.\quad}
\]
For every $\mu$, $G_c$ is the unique maximal submodule of $Q_c$, the
quotient $Q_c/G_c$ is simple, and $Q_c$ is indecomposable; for
$\mu\in\Z$ the sequence $0\to G_c\to Q_c\to Q_c/G_c\to0$ does not
split.
\end{theorem}
\begin{proof}
By \eqref{eq:FGXZ}, $M$ is simple. The PBW basis makes $F$ injective and
$H$ semisimple on $M$; $E=e_{-c}$ is locally nilpotent on $M$ because
$-c>a$; and $\ad E,\ad F$ are locally nilpotent on $U$: indeed
$[e_r,e_s]=[f_r,f_s]=0$ and, by \eqref{eq:brackets},
$[e_r,[e_r,[e_r,x]]]=[f_r,[f_r,[f_r,x]]]=0$ for every generator $x$, the
chains being $f_s\mapsto h_{r+s}\mapsto e_{2r+s}\mapsto0$ and the
analogous one for $f_r$. Thus Theorem~\ref{thm:containment} applies: every
proper submodule of $Q_c$ is contained in $G_c$. By Proposition~\ref{prop:G},
$G_c=0$ exactly when $\mu\notin\Z$, and $0\ne G_c\ne Q_c$ when
$\mu\in\Z$. The criterion and the structural assertions follow.
\end{proof}

\section{Central reduction and the derivation}
\label{sec:central}

\begin{proposition}[Reduction by commuting endomorphisms]
\label{prop:reduction}
Let $V$ be a $U$-module, let $U_F$ be an Ore localization flat as a right
$U$-module, and let a commutative algebra $R$ act on $V$ by $U$-module
endomorphisms. The action extends to $D_FV$ by $z(s\otimes v)=s\otimes zv$
and descends to $\T_FV$. For any algebraic ideal $I\subseteq R$,
\[
 D_F(V/IV)\cong D_FV/I(D_FV),\qquad
 \boxed{\quad\T_F(V/IV)\cong(\T_FV)/I(\T_FV),\quad}
\]
and injectivity of $F$ on $V/IV$ is not required.
\end{proposition}
\begin{proof}
The extension is well defined because each $z$ commutes with the
$U$-action, hence with $F^{-1}$: $F^{-m}(zv)=z(F^{-m}v)$. Localizing
$0\to IV\to V\to V/IV\to0$ and using flatness gives
$D_F(V/IV)\cong D_FV/\im D_F(IV)$; the image equals $I(D_FV)$, since
every element is a finite sum of terms $F^{-m}(zv)=z(F^{-m}v)$. The image
of $V/IV$ in $D_FV/I(D_FV)$ is $(\ell(V)+I(D_FV))/I(D_FV)$, whose cokernel
is $D_FV/(\ell(V)+I(D_FV))$; quotienting $\T_FV=D_FV/\ell(V)$ by
$I(\T_FV)$ produces exactly the same expression.
\end{proof}

At the critical level $\kappa=-2$ the Segal--Sugawara operators are the
normally ordered sums
\begin{equation}\label{eq:sugawara}
 T(n)=\frac14\sum_{j\in\Z}
 \bigl(2:e_{-j}f_{n+j}:+2:f_{-j}e_{n+j}:+:h_{-j}h_{n+j}:\bigr),
\end{equation}
where $:x_py_q:=x_py_q$ if $p<0$ and $:=y_qx_p$ if $p\ge0$. On a smooth
module each sum is finite, and at $\kappa=-2$ the $T(n)$ commute with
$\g$ \cite[Section 2.5]{FGXZ}.
\begin{lemma}\label{lem:sugawara}
Let $V$ be a smooth module of level $-2$ and $F=f_c$. The natural extension
of the $U$-endomorphism $T(n)$ from $V$ to $D_FV$ agrees with the actual
Sugawara action there; in particular $T(n)(F^{-m}v)=F^{-m}T(n)v$.
\end{lemma}
\begin{proof}
Localization preserves smoothness: for $m\ge1$,
\[
 h_rF^{-m}=F^{-m}h_r+2mF^{-m-1}f_{r+c},\qquad
 e_rF^{-m}=F^{-m}e_r-mF^{-m-1}(h_{r+c}+r\delta_{r+c,0}K)
 -m(m+1)F^{-m-2}f_{r+2c},
\]
and the $f_r$ commute with $F$; for a fixed PBW basis vector choose
$r,r+c,r+2c$ large enough that all resulting elements annihilate the
vector and the central $\delta$-term vanishes. The map $V\to D_FV$
intertwines $T(n)$ (each finite product intertwines, and normal ordering
makes the sum finite on each vector). Since $T(n)$ commutes with $\g$ it
commutes with $F$, hence with $F^{-1}$, which gives the displayed formula.
\end{proof}
\begin{corollary}
\label{cor:critical}
Let $V$ be a smooth module of level $-2$, fix integers $n_i$ and scalars
$\theta_i$ ($i\in J$), and let $R=\C[z_i:i\in J]$ act by $z_i=T(n_i)$ and
$I=(z_i-\theta_i:i\in J)$. For every $c\in\Z$,
\[
 \T_{f_c}(V/IV)\cong(\T_{f_c}V)/I(\T_{f_c}V),
\]
and the same holds for any algebraic ideal of $R$.
\end{corollary}
\begin{proof}
Apply Proposition~\ref{prop:reduction} and Lemma~\ref{lem:sugawara}.
\end{proof}

\section{The subspace \texorpdfstring{$G_c$}{Gc}: established facts and open questions}
\label{sec:Gc}

Assume $c\le-a-1$. Applying the automorphism $\sigma_{-c}$ of $\g$ defined
by $e_i\mapsto e_{i-c}$, $f_i\mapsto f_{i+c}$,
$h_i\mapsto h_i-c\delta_{i,0}K$, $K\mapsto K$, we may assume $c=0$, so
that $E=e_0$, $F=f_0$, $H=h_0$, and $Q_c$ is the quotient of
$M_{a+c,b-c}(\chi')$ at $f_0$, with $\chi'(h_0)=\mu$. We work in these
coordinates in this section.

Let $V_0\subset M$ be the span of the PBW basis vectors not involving
$f_0$, so $M=\bigoplus_{r\ge0}f_0^rV_0$. For $x\in V_0$ define linear
maps $A_r:V_0\to V_0$ by the finite expansion
\begin{equation}\label{eq:Ar}
 Ex=\sum_{r\ge0}f_0^rA_rx,\qquad A_0=\pi_0\circ\ad E,
\end{equation}
where $\pi_0$ projects onto $V_0$. The following replaces the incorrect
formula of an earlier draft; it was corrected after the independent audit
of 2026-09-22, whose counterexamples are reproduced in
\texttt{work/check\_audit\_regressions.py}.

\begin{lemma}[Corrected commutation formula for $E$]
\label{lem:Eaction}
If $x\in V_0$ with $Hx=(\mu+\delta)x$ and $m\ge1$, then in $Q_c$,
\begin{equation}\label{eq:Eaction}
 E(F^{-m}x)=\sum_{r=0}^{m-1}F^{-(m-r)}A_rx
 -m(\mu+\delta+m+1)F^{-m-1}x.
\end{equation}
\end{lemma}
\begin{proof}
By \eqref{eq:triple} and \eqref{eq:Ar},
\[
 EF^{-m}x=F^{-m}Ex-mF^{-m-1}Hx-m(m+1)F^{-m-2}Fx
 =\sum_{r\ge0}F^{-m}F^rA_rx-m(\mu+\delta)F^{-m-1}x
 -m(m+1)F^{-m-1}x;
\]
the terms with $r<m$ survive as $F^{-(m-r)}A_rx$, and the terms with
$r\ge m$ lie in $M$ and vanish in $Q_c$.
\end{proof}

\begin{lemma}\label{lem:top}
If $0\ne v\in Q_\nu$ and $Ev=0$, then $\nu\in\Z_{\ge0}$, and in the
expansion $v=\sum_{m=1}^{N}F^{-m}x_m$ with $x_m\in V_0$ and $x_N\ne0$ one
has $N=\nu+1$.
\end{lemma}
\begin{proof}
The coefficient of $F^{-(N+1)}$ in $Ev$ receives contributions only from
the diagonal term $-N(\mu+\delta_N+N+1)x_N$ of \eqref{eq:Eaction}, where
$\mu+\delta_N=\nu-2N$; the sum $\sum_rF^{-(m-r)}A_rx_m$ has negative
powers at most $m\le N$. Thus $0=-N(\nu-N+1)x_N$ forces $\nu=N-1$.
\end{proof}

\begin{proposition}[Triangular system for the vectors annihilated by $E$]
\label{prop:recursion}
A nonzero vector $v=\sum_{m=1}^{\nu+1}F^{-m}x_m\in Q_\nu$ satisfies
$Ev=0$ if and only if
\begin{equation}\label{eq:recursion}
 \sum_{m=k}^{\nu+1}A_{m-k}x_m=(k-1)(\nu-k+2)x_{k-1},
 \qquad 1\le k\le\nu+1,\quad x_0=0.
\end{equation}
For $\nu=1$ this reads $x_1=A_0x_2$ and $(A_0^2+A_1)x_2=0$. The earlier
description of the primitive space by the single condition
$A_0^{\nu+1}x_{\nu+1}=0$ is false: the element
$Z=h_{-1}^2+4f_{-1}e_{-1}+2h_{-2}$ commutes with $E,F,H$, and at
$\mu=-1$ the vector $p=\lambda_1F^{-1}u+F^{-2}f_1u$ lies in $\ker E$
(weight $1$) while $Zp$ has layers $x_1=(\lambda_1Z-4h_{-1})u$,
$x_2=(f_1Z-4(\kappa+1)f_{-1})u$ with $A_0^2x_2=8e_{-1}u\ne0$ and
$A_1x_2=-8e_{-1}u$, so that \eqref{eq:recursion} holds with $k=1$ by
cancellation.
\end{proposition}
\begin{proof}
Equate the coefficient of $F^{-k}$ in $Ev=0$ using \eqref{eq:Eaction}:
from $\sum_rF^{-(m-r)}A_rx_m$ with $m-r=k$ (i.e.\ $r=m-k\ge0$) we get
$\sum_{m\ge k}A_{m-k}x_m$, and the diagonal contributes
$-(k-1)(\mu+\delta_{k-1}+k)x_{k-1}=-(k-1)(\nu-k+2)x_{k-1}$, using
$\mu+\delta_{k-1}=\nu-2(k-1)$. The assertions about $Z$ are direct
bracket computations ($[E,Z]=[F,Z]=[H,Z]=0$; the displayed values are
verified in \texttt{work/check\_audit\_regressions.py}).
\end{proof}

\begin{proposition}[Decomposition into irreducibles of the subalgebra generated by $E,H,F$]
\label{prop:sl2}
Let $\mathfrak s$ be the subalgebra generated by $E,H,F$, isomorphic to
$\slc$. As an $\mathfrak s$-module,
\[
 G_c\cong\bigoplus_{\nu\in\Z_{\ge0}}P_\nu\otimes L(\nu),
\]
where $L(\nu)$ is the $(\nu+1)$-dimensional irreducible module and
$P_\nu=\ker E\cap G_c\cap(Q_c)_\nu$ is the multiplicity space (possibly
infinite-dimensional). This is not a decomposition into
$\g$-submodules.
\end{proposition}
\begin{proof}
$E$ and $F$ are locally nilpotent on $G_c$, so every vector lies in the
finite span $\Span\{F^rE^sv\}$; complete reducibility of the finite
dimensional modules of $\slc$ gives the decomposition. A vector
$p\in\ker E$ with $F^Np=0$ generates $L(\nu)$ with $\nu$ its $H$-weight,
and $F^Np=0$ forces $\nu\le N-1$, so $\nu\in\Z_{\ge0}$.
\end{proof}

\begin{proposition}[A projection lowering the highest weight by $2$]
\label{prop:lowering}
For $\nu\ge2$, $p\in P_\nu$ and $j\in\Z$,
\begin{equation}\label{eq:lowering}
 \ell_j(p):=f_jp-\tfrac1\nu Fh_jp
 -\tfrac{1}{\nu(\nu+1)}F^2e_jp\in P_{\nu-2}.
\end{equation}
It is not identically zero: at $\mu=-2$, for $p_2=e_{-1}w_1\in P_2$,
\[
 \ell_1(p_2)=F^{-1}\Bigl(\tfrac13f_1e_{-1}+\tfrac{\lambda_1}{6}h_{-1}
 +\tfrac{2\kappa-2}{3}\Bigr)u+M\ne0.
\]
\end{proposition}
\begin{proof}
Compute $E\ell_j(p)$ term by term. $Ef_jp=[E,f_j]p+f_jEp=h_jp$. Since
$Hh_jp=\nu\,h_jp$ and $Eh_jp=[E,h_j]p+h_jEp=-2e_jp$,
\[
 EFh_jp=[E,F]h_jp+FEh_jp=Hh_jp-2Fe_jp=\nu\,h_jp-2Fe_jp.
\]
Since $e_jp\in P_{\nu+2}$ (as $[E,e_j]=0$), we have $EFe_jp=(\nu+2)e_jp$
and
\[
 EF^2e_jp=[E,F]Fe_jp+FEFe_jp=HFe_jp+(\nu+2)Fe_jp=2(\nu+1)Fe_jp,
\]
because $HFe_jp=FHe_jp-2Fe_jp=(\nu+2-2)Fe_jp=\nu Fe_jp$. Hence
\[
 E\ell_j(p)=h_jp-\tfrac1\nu(\nu h_jp-2Fe_jp)
 -\tfrac{1}{\nu(\nu+1)}2(\nu+1)Fe_jp=0,
\]
and all three terms of \eqref{eq:lowering} have $H$-weight $\nu-2$. The
displayed value of $\ell_1(p_2)$ is obtained by the direct PBW computation
reproduced in \texttt{work/check\_audit\_regressions.py}.
\end{proof}

\begin{remark}[Withdrawn claims and what is open]
\label{rem:withdrawn}
An earlier draft claimed that the subspaces
$G^{\ge\nu}:=\bigoplus_{\nu'\ge\nu}P_{\nu'}\otimes L(\nu')$ are
$\g$-submodules and deduced that $G_c$ is never simple, that
$\soc(G_c)=\soc(Q_c)=0$, and that $Q_c$ has infinite length. These claims
are withdrawn. Counterexample: at $\mu=-2$ and $\kappa\ne0$, for
$p=w_1\in L(0)$ we have $h_{\pm1}p=-Fe_{\pm1}p\in G^{\ge2}$ (indeed
$h_rp=e_rFp-Fe_rp$ and $Fp=0$, while $e_{\pm1}p\in P_2$); if $G^{\ge2}$
were a submodule, $[h_1,h_{-1}]p=2\kappa p\in G^{\ge2}$, contradicting
$p\in L(0)$. Moreover the projection of Proposition~\ref{prop:lowering}
maps $P_2$ into $P_0$, so these subspaces are not stable even at
$\kappa=0$. In the case $\mu\in\Z$ the containment theorem still gives
$\soc(Q_c)=\soc(G_c)$ (every simple submodule of $Q_c$ is proper and hence
contained in $G_c$), but the value of this socle, the simplicity of $G_c$,
and the composition length of $Q_c$ are open.
\end{remark}

\section{Open problems}
\label{sec:open}
\begin{enumerate}[leftmargin=*,label=\arabic*.]
\item \textbf{Simplicity in the case $-a\le c<b$.} Here the element
$e_{-c}$ is injective on $Q_c$ (it is a PBW basis element of the
complement of $S_{a,b}$; grading by the weight $2\#e+\#h$, its action on
the leading term is multiplication by an independent variable, and the
commutation corrections have lower weight), so $G_c=0$; moreover
$Q_c=Uw_1$ when $\lambda_L\ne0$, by $[e_{L-c},f_c]=h_L$ and
$e_{L-c}w_m=-m\lambda_Lw_{m+1}$. The proof of Theorem~\ref{thm:deep} does
not apply because $e_{-c}$ is not locally nilpotent on $M$. The simplicity
criterion for $Q_c$ is unknown.
\item \textbf{Submodules of $G_c$.} For $\mu\in\Z$, determine whether
$G_c$ is simple, compute $\soc(G_c)=\soc(Q_c)$, and decide whether $Q_c$
has finite composition length; a correct analysis must track the three
channels $P_\nu\to P_{\nu+2},P_\nu,P_{\nu-2}$ simultaneously
(Propositions~\ref{prop:recursion} and \ref{prop:lowering}).
\item \textbf{Critical reduction.} Corollary~\ref{cor:critical} gives
$\T_{f_c}(V/IV)\cong(\T_{f_c}V)/I(\T_{f_c}V)$, but the simplicity of the
reduced quotients at $\kappa=-2$ is not determined.
\item \textbf{Parameter classification.} A complete isomorphism
classification of the $Q_c$, including $\lambda_L=0$ and $\kappa=-2$, is
open.
\end{enumerate}

\begin{thebibliography}{9}
\bibitem{FGXZ}
V.~Futorny, X.~Guo, Y.~Xue and K.~Zhao,
\emph{Smooth representations of affine Kac--Moody algebras},
arXiv:2404.03855v2 (2024).
\url{https://arxiv.org/abs/2404.03855v2}.

\bibitem{Unified}
Project note, \emph{Boundary and Nonboundary Localization Quotients of
Smooth Affine Modules}, September 2026 (unified note and its addendum of
September 15, 2026).

\bibitem{Audit}
Independent audit of 2026-09-22,
\texttt{audits/2026-09-22/audit\_report\_cn.md} and
\texttt{audits/2026-09-22/independent\_pbw\_audit.py} (154 exact checks);
counterexamples reproduced in
\texttt{work/check\_audit\_regressions.py} (24 checks).
\end{thebibliography}

\end{document}